\documentclass[12pt]{article}
\usepackage[T1]{fontenc}
\usepackage[utf8]{inputenc}
\usepackage{lmodern}
\usepackage{textcomp}
\usepackage{enumitem}
\usepackage{geometry}
\usepackage{graphicx}
\usepackage{amsmath, amssymb}
\geometry{margin=1in}
\begin{document}
\begin{center}
Physics - Undergraduate Level Assessment 12 \
Total Marks: 40 \
Answer all questions.
\end{center}

\section*{Multiple Choice (10 marks)}
\begin{enumerate}[label=\arabic*.]
\item A meter stick with a speed of 0.8 c moves past an observer. In the observer's reference frame, how long does it take the stick to pass the observer ?
\begin{enumerate}[label=(\alph*)]
    \item 1.6 ns
    \item 2.5 ns
    \item 4.2 ns
    \item 6.9 ns
\end{enumerate}
\item A grating spectrometer can just barely resolve two wavelengths of 500 nm and 502 nm, respectively. Which of the following gives the resolving power of the spectrometer?
\begin{enumerate}[label=(\alph*)]
    \item 2
    \item 250
    \item 5,000
    \item 10,000
\end{enumerate}
\item The sign of the charge carriers in a doped semiconductor can be deduced by measuring which of the following properties?
\begin{enumerate}[label=(\alph*)]
    \item Specific heat
    \item Thermal conductivity
    \item Electrical resistivity
    \item Hall coefficient
\end{enumerate}
\item The speed of light inside of a nonmagnetic dielectric material with a dielectric constant of 4.0 is
\begin{enumerate}[label=(\alph*)]
    \item 1.2 * $10^9$ m/s
    \item 3.0 * $10^8$ m/s
    \item 1.5 * $10^8$ m/s
    \item 1.0 * $10^8$ m/s
\end{enumerate}
\item The negative muon, mu^-, has properties most similar to which of the following?
\begin{enumerate}[label=(\alph*)]
    \item Electron
    \item Meson
    \item Photon
    \item Boson
\end{enumerate}
\end{enumerate}

\section*{True / False (10 marks)}
\textit{Answer True or False}
\begin{enumerate}[label=\arabic*.]
\setcounter{enumi}{5}
\item True or False: Electromagnetic radiation emitted from a nucleus is most likely in the form of gamma rays.
\item True or False: The primary source of the Sun's energy is the fusion of four hydrogen atoms into one helium atom, releasing energy according to E=$mc^2$.
\item True or False: The resolving power of the grating spectrometer for wavelengths of 500 nm and 502 nm is 250.
\item True or False: A dye laser is the best type of laser for spectroscopy across a range of visible wavelengths.
\item True or False: In a single-electron atom with l = 2, there are only 3 allowed values for the magnetic quantum number m\_l.
\end{enumerate}

\section*{Long Form Response (20 marks)}
\textit{Answer each question with a clear and complete explanation.}
\begin{enumerate}[label=\arabic*.]
\setcounter{enumi}{10}
\item Discuss the allowed transitions for an electron in the n = 4, l = 1 state of hydrogen. Explain which final states are not reachable and why, using quantum mechanical principles.
\item Analyze a one-dimensional inelastic collision where a particle of mass 2 m collides with a stationary mass m. Discuss the initial kinetic energy and calculate the fraction lost during the collision.
\end{enumerate}

\vfill
\noindent\textit{End of Paper}
\end{document}